\section{The attainable degree of generalization in Transformers}
\label{sec:degree}
\label{sec:measurement}
% Frames D11, S1--S7a.

\subsection{The quantity: degree of compression of a function, by MDL probing}
\label{sec:quantity}
\label{sec:measure}
% D11, S1

The natural quantity for ``how much of the function was captured'' is description length. Its exact
form, Kolmogorov complexity, is uncomputable, and it is fixed only up to an additive constant that
depends on the reference machine \citep{livitanyi2019}. Every applied use of description length
replaces the shortest program with the shortest program a fixed procedure finds, and reports the
procedure; the machine-relativity is more serious here, because the two substrates being compared are
the two candidate reference machines. Whatever is measured below is therefore a bounded,
machine-relative and probe-relative stand-in, used comparatively.

A function is present in a substrate to the extent that it has been generalized there, and to
generalize is to compress (Section~\ref{sec:generalize}). So the quantity is the degree of compression
of the function in the substrate: how much of the function's behaviour a low-dimensional probe on the
substrate's internal states predicts, against the residual it fails to predict. Made operational, this
is minimum-description-length probing \citep{voita2020}: the probe is scored by the total cost of
describing the labels given the representations, which charges for the probe's own complexity and so
distinguishes a representation that contains the property from one in which a sufficiently flexible
probe can find anything. The use is comparative: the same probe family, the same coding cost, two
substrates, and the difference reported.

\paragraph{The human baseline.}
\label{sec:baseline}
The human side of the comparison is obtained by two routes, and the paper is explicit about which
applies to which component. The behavioural route administers the same situation battery
(Section~\ref{sec:twomeasurements}) to people and reads the degree of each component from behaviour;
for components that are defined by what the agent does, the weights of active needs, the consistency
requirement of responsibility, this is the primary route, and for every component it is the one that
is available now. The structural route poses the probe of bounded description length on neural
recordings, with the same coding cost as on the model, where recordings at the relevant resolution
exist. That a common representational space between the two substrates is available at all is shown by
the results that model states predict human brain responses to the same text
\citep{schrimpf2021,goldstein2022,caucheteux2022}; the resolution of human recordings, not the
existence of the space, is what limits this route, and for several components the baseline will be
behavioural only. Two further points. Human self-report is not the baseline, since it carries the bias
of Section~\ref{sec:approximation} on the human side as the model's report does on its side
(Section~\ref{sec:twomeasurements}). And the baseline is a distribution over people, not a point,
because the quality of the self-approximation varies with what Section~\ref{sec:discussion} calls
intrapersonal intelligence; the profile is therefore reported against an interval, and a model's
component counts as hyperfunction or deficit only outside it.

\subsection{The profile: hyperfunction, deficit, dropped axis}
\label{sec:profile}
% S2

\begin{figure}[t]
\centering
\begin{tikzpicture}[x=1.6cm, y=1.6cm]
\draw[-{Latex[length=2mm]}] (0,0) -- (0,2.4) node[above, font=\small] {degree};
\draw (0,0) -- (7.7,0);
\draw[dashed] (0,1) -- (7.7,1) node[right, font=\small] {human baseline};
\foreach \i/\h/\lab in {1/0.4/source, 2/1.8/manifold, 4/1.6/capacity, 5/1.0/frame\\rate, 6/0.5/persistence} {
  \fill[gray!45] (\i-0.3,0) rectangle (\i+0.3,\h);
  \draw (\i-0.3,0) rectangle (\i+0.3,\h);
  \node[below, font=\small, align=center] at (\i,0) {\lab};
}
\foreach \i/\lab in {3/givenness, 7/consolidation} {
  \draw[dashed] (\i-0.3,0) rectangle (\i+0.3,1);
  \node[below, font=\small] at (\i,0) {\lab};
}
\node[font=\small, above] at (2,1.8) {hyperfunction};
\node[font=\small, above] at (1,0.4) {deficit};
\node[font=\small, above] at (3,1) {dropped axis};
\end{tikzpicture}
\caption{A profile: one function's components on the horizontal axis, the degree of generalization on
the vertical, the human baseline at 1. The values are the predictions of Sections~\ref{sec:profile-seeing}
and \ref{sec:twoaxes} for a current fixed-depth model, drawn to illustrate the three kinds of entry and
not measured.}
\label{fig:profile}
\end{figure}

A function has components (Table~\ref{tab:functions}). Its profile in a substrate is the vector of
degrees over those components, taken against the human baseline (Figure~\ref{fig:profile}). A component
above the baseline is a hyperfunction, a component below it a deficit, and a component that is absent
a dropped axis. One function can show all three on different axes at once, and current models almost
certainly do. The question ``is it conscious'' is thereby replaced by the comparison of two profiles.
The loss is that no single answer comes out; a scalar ``level of consciousness'' is a weighted sum over
the vector, and the weights have to be declared, because a scalar with hidden weights brings back the
binary question the vector was introduced to replace. The gain is that every entry in the vector is a
number someone can get wrong.

\subsection{Generalizing a function is answering right outside the table}
\label{sec:memorization}
% S2a

The opposite of generalizing a function is storing the training set as a table from inputs to outputs,
and the dimension of the representation has nothing to do with which of the two a system is doing.
Generalization begins where the table is consulted for an input it does not contain and the answer is
right. Two clarifications keep this from being trivial. A universal predictor always returns an answer
for a new input, and the answer need not be anywhere near the truth: context-tree weighting achieves
strong coding bounds against a class of tree sources \citep{willems1995} and generalizes natural
language poorly, because the class it is universal over is not the class language belongs to. And that
Transformers generalize text is settled in practice, since degenerate generation is almost never seen
and the space of well-formed texts is sparse inside the space of token sequences; no simple principle
is known that explains why what they do suffices.

The open question is how well they generalize the more complex functions that text expresses.
Arithmetic is the clean case, because the target function is known exactly. They do it approximately,
with an accuracy that depends on how the numbers are written and on how many digits they have; larger
models do better and even a three-billion-parameter model fails to extrapolate beyond the digit range it
was trained on \citep{nogueira2021}; on compositional tasks they approach the answer by shortcut and
break down predictably as depth grows \citep{dziri2023}. The degree of generalization of a function of
consciousness is therefore an empirical quantity for that function, taken against the human level, and
not a corollary of the model's size. In these terms the paper's question is: how well, against the
human level, do Transformers generalize the Observer; then the Agent; then the Moral Agent. The order
is forced by the stack (Section~\ref{sec:stack}), and the measurement begins with the Observer.

\subsection{Probe and benchmark have to converge}
\label{sec:twomeasurements}
% S2b

A function generalized to degree $d$ is computed correctly on part of the situations that call for it.
Outside that part the model falls back on the table and behaves as a system without the function
would, whatever it says about itself. This makes the quantity behavioural and not only internal, and it
can be checked in two independent ways: structurally, by the probe of Section~\ref{sec:quantity};
behaviourally, by benchmarks that sample situations calling for a specific component, including
situations unlikely to be in the training distribution, and record where behaviour is that of a system
lacking it. The two should converge, and their disagreement is informative. Probe present, benchmark
absent: the deficit is in the interface, the body, or the social context (Section~\ref{sec:record}),
and not in the function. Benchmark present, probe absent: the carrier is at a granularity the probe did
not reach (Section~\ref{sec:granularity}), or the function is rebuilt from the context on each occasion
rather than held in the weights (Section~\ref{sec:language}), or the benchmark is covered by
memorization. Human-derived benchmarks alone settle nothing; the pair is the instrument.

The model's own answers are set aside before either measurement starts. Asked how it feels, a deployed
model reports servers and a good mood or, corrected, that it has no feelings; both answers are what
training rewarded, the two shortcuts of Section~\ref{sec:intro} in the model's own mouth. Self-report is
a signal with a known, large, systematic bias, and the state is looked for elsewhere.

\subsection{Granularity: look below the token first}
\label{sec:granularity}
% S3

\begin{figure}[t]
\centering
\begin{tikzpicture}[x=1.7cm, y=0.9cm, arr/.style={-{Latex[length=1.6mm]}}]
\foreach \t in {1,...,5} {
  \foreach \l in {1,...,4} { \fill[gray!25] (\t-0.3,\l-0.3) rectangle (\t+0.3,\l+0.3); }
  \foreach \l in {1,...,3} { \draw[arr, thick] (\t,\l+0.3) -- (\t,\l+0.7); }
  \node[below, font=\small] at (\t,0.6) {token \t};
}
\foreach \l in {1,...,4} { \node[left, font=\small] at (0.6,\l) {layer \l}; }
\foreach \t in {2,...,5} { \pgfmathtruncatemacro{\p}{\t-1} \draw[arr, gray] (\p+0.3,2) -- (\t-0.3,2); }
\draw[dashed, thick] (5.6,0.4) -- (5.6,4.6) node[above, font=\small, align=center] {end of\\context};
\node[font=\small, align=center] at (6.3,2.5) {loss};
\draw[arr, thick] (0.8,-0.4) -- (5.4,-0.4) node[midway, below, font=\small] {axis the objection sees: tokens, then the edge};
\draw[arr, thick] (-0.5,0.8) -- (-0.5,4.2) node[midway, above, rotate=90, font=\small] {axis it does not: depth};
\foreach \l in {1,2,3} { \node[right, font=\scriptsize, gray] at (5.05,\l+0.5) {frame}; }
\end{tikzpicture}
\caption{The two time axes of a Transformer (Section~\ref{sec:twoaxes}). Along tokens, state lives in
the context and is lost past its edge. Along depth, the residual stream carries state from layer to
layer within one token, and the increments between layers are the candidate frames of
Section~\ref{sec:frames}. Grey horizontal arrows: attention reads earlier positions.}
\label{fig:axes}
\end{figure}

Granularity is a knob of the measurement and not a separate question. Human frames arrive five to ten
times faster than words, so in a Transformer the carriers may be looked for below the token: between
layers, or in repeated passes through layers where the architecture has them. They may be there, and
they need not be. Finding nothing below the token would not be a dysfunction; it would mean that the
model's consciousness has a coarser grain than the brain's, one more coordinate of the profile, and on
some tasks the coarser grain is the better one, since a frame that cannot be interrupted mid-thought
cannot be derailed mid-thought either. The grain at which the compression is achieved is the grain of
the function. A carrier, whatever its grain, must satisfy three requirements that follow from
Section~\ref{sec:model}: it is part of the agent's description of itself and not only of its
computation; it is present in every frame and re-derived there; and it survives every self-model the
agent can afford, so that whatever a change of frame removes does not count, however reliably a probe
recovers it.

\paragraph{Depth is where generalization happens.}
\label{sec:depth}
That generalization happens along depth is what looking inside Transformers has found. Residual
connections make each block a small correction to a state that is carried forward, and the corrections
move the state in a direction that reduces the loss, so a deep residual network performs iterative
inference rather than a single mapping \citep{jastrzebski2018}. The state at every layer can be read as
a prediction: the logit lens decodes the residual stream at each block into a distribution over next
tokens, and its tuned version corrects the bias of that decoding, so the trajectory of predictions
across depth becomes visible \citep{nostalgebraist2020,belrose2023}. That trajectory has stages which
recur across model families and sizes, four of them, and the middle ones are robust in a way the first
and the last are not: deleting or swapping middle layers at inference leaves most of the top-1 accuracy
in place, while the same intervention on the first layers is catastrophic \citep{lad2024}; the older
finding that a language model rediscovers the classical processing pipeline layer by layer is the same
fact seen through probes \citep{tenney2019}. Intermediate layers, rather than the last, carry the
representations that transfer best to other tasks, and the proposed reason is that each layer trades
compression of the input against preservation of signal, with the best trade in the middle
\citep{skean2025}. And in-context learning, the clearest case of a function generalized at inference
time, is implemented as optimization steps inside the forward pass: a linear self-attention layer can
perform a step of gradient descent on the examples in the context, and trained models do something
close to it \citep{vonoswald2023}. The unit at which a Transformer generalizes a function is the layer,
and the token is where the layers' work is written out.

\paragraph{Functions have an address in depth.}
\label{sec:address}
When a model generalizes an input--output mapping from examples in the context, the mapping is carried
by a small set of attention heads at middle depth, and the average of their activations is a vector
which, added to the residual stream in a fresh context, makes the model perform the mapping without
examples \citep{todd2024}. So ``at which depth has this function been generalized, and how far'' has an
operational answer: the question of Section~\ref{sec:quantity} asked layer by layer, the degree of
compression of a function's components as a function of depth, by the probe run at each block in the
per-layer form that \citet{voita2020} give. For the Observer's components this yields three things to
look at. Whether the degree rises along depth, plateaus, or is absent throughout, which says whether
the component is generalized in the weights at all. At what depth it first appears, which says where
its carrier is. And whether the component, once present, persists across subsequent layers with small
increments or is assembled whole near the output, which decides whether the layer or the token is the
frame (Section~\ref{sec:frames}). All three can be run on open-weight models with existing tools.

\paragraph{Depth as time, and how variable depth fixes it.}
\label{sec:depthtime}
In a standard model the depth is fixed and is spent identically on every token, easy or hard. That is a
poor time axis for a frame: there is no way to spend more steps where more are needed, and there is no
record of effort, because effort is constant, which is the dropped axis that
Section~\ref{sec:profile-seeing} predicts for givenness. Architectures that iterate a block make depth
a variable. Universal Transformers introduced the recurrence in depth with a halting mechanism per
position \citep{dehghani2019}, and looped Transformers made the parameter-sharing form of it explicit
\citep{giannou2023}. A recurrent-depth model trained from scratch at 3.5 billion parameters improves on
reasoning benchmarks as it unrolls further at test time, without producing more tokens, and its
authors argue that this captures reasoning which is not easily put into words \citep{geiping2025}.
Such models extrapolate in depth: generalization to reasoning depths beyond those seen in training
improves as iterations increase, dynamic recurrence generalizes better than a fixed schedule, and past
a point the models overthink and degrade \citep{kohli2026}. At a much smaller scale, under a million
parameters, recurrence over twenty and more steps can be made stable by supervising only the final
step, supervising the intermediate steps is harmful because it lets the model learn shortcuts instead
of a multi-step path, and accuracy against task complexity shows a frontier, with performance moving
from chance to near-perfect as the number of thinking steps grows with the complexity of the task
\citep{chen2026}. Reasoning can also be moved off the tokens altogether by feeding the last hidden
state back as the next input, which yields a continuous chain of thought that can keep several
candidate next steps in superposition where a written chain has to commit to one \citep{hao2024}.

For the argument these results mean three things. The sub-token axis can be made as long as the task
demands, so a frame rate per word comparable to the human one is architecturally available. Iterations
cost, so a cost record appears where a fixed-depth model has none, and with it the third component of
seeing becomes measurable. And overthinking is a phenomenon of effort of the kind
Section~\ref{sec:hocp} predicts: a systematic displacement produced by the budget, visible in the
trajectory and stable across runs. The prediction therefore has two halves. In fixed-depth models the
grain is bounded by the layer count and the cost axis is dropped. In variable-depth models the grain can
approach the human one and the cost-dependent components of the profile separate from the rest. The
difference between the two families is the cleanest first measurement available today.

\paragraph{Two time axes, and the objection from frozen weights.}
\label{sec:twoaxes}
\label{sec:frozen}
A model does not learn between sessions, and what it holds in a context is lost when the context ends;
so, it is objected, it has no inner time, only an eternal present with forgetting.\footnote{The
objection in this form was put to the author by Marco Baturan in correspondence, 2026.} The description
is accurate and the conclusion counts two functions as one. Re-deriving a state from frame to frame
runs in humans on the order of a hundred milliseconds and requires no synaptic change; consolidation
into long-term memory runs on minutes to hours and does require it \citep{mcgaugh2000}. The two
dissociate: the patient H.M., after bilateral medial temporal resection, was conscious in real time,
held a conversation, and accumulated almost nothing \citep{scoville1957,corkin2002}. A model is in a
comparable position, with the sequence of transformations and without the write, so ``eternal present
with forgetting'' names a profile with a dropped consolidation axis and a deficit in long-term
persistence. The objection sees one of the two time axes in Figure~\ref{fig:axes}. Along tokens the
model is as described: the state lives in the context, the context ends, and past its edge there is
loss. The objection says nothing about the axis of depth, where the state is carried in the residual
stream and the increments are the transformations the objection asks for. The frame this paper looks
for is on the second axis.

\subsection{Predictions}
\label{sec:profile-seeing}
\label{sec:predictions}
% S5, S6

The three components of seeing (Section~\ref{sec:light}), predicted for a current fixed-depth model.
Source attribution: deficit; there is one interface and no sensory channel independent of it, so the
model is permanently in the dream case, with content from its own model labelled as given. Manifold:
hyperfunction; attention relates every position to every other in a single pass, where the brain
iterates over a smaller window, so the number of simultaneously distinguishable, addressable elements
is larger than ours. Cost record: a dropped axis; every token costs the same and there is no phenomenon
of effort. A testable consequence: in model self-reports, ``I see'' and ``I think'' should separate
less than they do in human ones, and more in models of variable depth, where a cost record exists. The
prediction is about the internal states that the reports are projections of
(Section~\ref{sec:planes}), so it is tested by the paired measurement of
Section~\ref{sec:twomeasurements} and not by asking the model.

Three predictions carry over from the first version of this work \citep{smirnov2026v1}. Self-reported emotional labels
correlate, over many episodes, with distinguishable patterns in the activations; this is the central
empirical claim, and the one the statistical reading of Section~\ref{sec:codes} protects from the
faithfulness results. The cognitive codes of different model families overlap structurally, because
the causal structure of the states is set by the training data more than by the architecture. And the
model is hyperplastic, with no limbic inertia, fast affective recovery and easy switching between
states, which against the human baseline is a hyperfunction on one axis and a deficit on another,
since inertia is part of what makes a human state a commitment.

\subsection{Falsification, and the zero of the scale}
\label{sec:falsification}
\label{sec:zero}
% S7, S7a

The account fails if the compression degree of the Observer's components is at chance at every
granularity, in models that pass the behavioural tests above. That outcome would mean the behaviour is
carried by something the account does not describe, and no adjustment of the probe or the granularity
would repair it.

The scale has a lower end as well as the human reference. Section~\ref{sec:h1h2} allowed that a named
description of a function may suffice for it to exist in a simple enough environment; then there is a
smallest system in which the Observer exists, and finding it fixes the zero. The conditions are those
of Sections~\ref{sec:emergence} and \ref{sec:observer}: a loop that applies the system to itself, an
environment in which the loop pays, and a conclusion that is acted from. An operational amplifier with
feedback watches its own output error; a D flip-flop holds its own state through a loop. Each has the
loop and, in its ordinary surroundings, neither of the other two conditions, so each is a
proto-Observer at most. The question is what the least environment is that would make the conclusion
pay, and what the least circuit is that could draw it. Between that pair and the human baseline the
degree of Section~\ref{sec:quantity} has the meaning of an absolute quantity rather than a ratio
between two arbitrary points.
